\documentclass[11pt]{article}
\usepackage[margin=1in]{geometry}
%%%%%%%%%%%%%%%%%%%%%%%%%%%%%%%%%%%%%%%%%%%%%%%%%%%%%%%%%%%%%%%%%%%%%%%%%%%%
\usepackage{amsmath,amssymb,amsfonts}
\usepackage{amsthm,latexsym,mathrsfs}
\usepackage{bm,bbm,dsfont}
\usepackage{color,graphicx}
\usepackage[latin1]{inputenc}
\usepackage[english]{babel}
\usepackage{xr-hyper}
\usepackage[colorlinks=true,linkcolor=blue,citecolor=blue,urlcolor=blue]{hyperref}
% labels of the main document, prefixed M-: build it first, so that its .aux exists
\externaldocument[M-]{math_bio_positive_matrices_clean}

\theoremstyle{plain}
\newtheorem{theorem}{Theorem}[section]
\theoremstyle{remark}
\newtheorem{remark}[theorem]{Remark}

\everymath{\displaystyle}
\numberwithin{equation}{section}
\renewcommand{\thesection}{S\arabic{section}}

\title{Supplement: numerical tests\thanks{Parts of this manuscript were edited with the assistance of an AI tool (Claude Opus 5.5, Anthropic). The authors are solely responsible for all proofs and statements.}}
\author{Zhilan \and John \and Yuliya \and Ludmil}
\date{}

\begin{document}

\maketitle

\begin{abstract}
Supplement to ``Minimizing the spectral radius of a diagonally scaled positive
matrix''. Equation, theorem and section numbers without a prefix S refer to
that document.
\end{abstract}

\section{Bounds for the single dose problem}\label{s-bounds}

The fixed term in~\eqref{M-split} confines the objective function: since
$1-\varepsilon_1\le g_p\le 1$ we have
$(1-\varepsilon_1)A_1\le A_1\operatorname{diag}(\bm g)\le A_1$, so by
Theorem~\ref{M-mono} in~\S\ref{M-sec-facts}
\begin{equation}\label{s-two-sided}
(1-\varepsilon_1)\,\rho(A_1)\le\rho\left(A_1\operatorname{diag}(\bm g)\right)\le\rho(A_1),
\end{equation}
with equality on the right exactly when no dose is allocated, $\chi_1=0$, and
the left approached as every $t_p\to\infty$. For the data of~\cite{JZh} the
interval is $\left\lbrack 0.8846136218,\ 4.4230681092\right\rbrack$ and the
constrained minimum of Section~\ref{M-sec-single} is $2.9059382878$, $43$ per cent
of the way down.

\section{A numerical test for single dose}\label{s-test}

For the data of~\cite{JZh}, with $m=n=4$, $N=16$ and $A_c$ of order $32$, the
right-hand side of~\eqref{M-constraints} is the weekly budget $b_{1i}$, the
number of doses of the primary series given in location $i$, and
$\sigma_{1,ij}$ is the number of susceptible individuals. The minimum
of~\eqref{M-t-program}, computed by the algorithm of
Section~\ref{M-sec-algorithm}, is $\rho=2.9059382878$, and the spectral radius of
the $32\times 32$ matrix $\operatorname{NGM}$ at the computed minimizer has the
same value. At the computed point all four constraints are active to a
relative $4.9\cdot 10^{-12}$, the ratio in~\eqref{M-kkt} has a spread of at most
$2.6\cdot 10^{-9}$ over the coordinates with $t_p>1$, the fourth coordinate,
the age group $60+$, is at $t_p=1$ for every $i$, and eight random feasible
starting points give values with a spread of $1.9\cdot 10^{-9}$ and minimizers
that agree to $8.3\cdot 10^{-11}$. The independent formulation of
Section~\ref{s-gp} agrees to $6.4\cdot 10^{-13}$.

An allocation of the same budget computed earlier and stored with the data,
with the same $\sigma_{1,ij}$ and the same $\operatorname{NGM}$, gives
$\rho=2.9247935891$; the minimum of Section~\ref{M-sec-single} is smaller by
$0.64$ per cent.

\section{A test instance with three doses}\label{s-three}

No budget for the second and third doses is given with the data
of~\cite{JZh}. On an instance built from these data, with the coefficients
$\sigma_{2,ij}$ and $\sigma_{3,ij}$ both taken equal to the numbers of partially
protected individuals, the daily budgets $b_{1i}/7$ for the first dose and
$b_{2i}=b_{3i}=b_{1i}/28$, the objective function of~\eqref{M-three-program} is
not convex in $\widehat{\bm u}$: $24$ of $4000$ random segments of the
product of simplices violate the convexity inequality, the worst by
$1.1\cdot 10^{-3}$, and twelve random starts of a local method return eight
distinct values between $3.8757$ and $3.8879$. With the weekly budgets $b_{1i}$
and $b_{2i}=b_{3i}=b_{1i}/4$ the same two tests find no violation in $4000$
segments, and the twelve starts end at one value, $1.661826$. With a single
dose and the weekly budget the segment test finds no violation in $4000$
segments, as Theorem~\ref{M-global} requires.

\section{The geometric programming form}\label{s-gp}

The program~\eqref{M-t-program} can be traded for one with a linear objective
function and no eigenvalue computation, which is the independent check quoted
in Section~\ref{s-test}. For $M\ge 0$ irreducible, with eigenvectors
$\bm v,\bm w>0$ of $M$ and $M^T$ corresponding to the spectral radius,
\begin{equation}\label{s-cw}
\rho(M)=\min\left\{\theta>0\ :\ \exists\,\bm x>0\ \mbox{with}\
M\bm x\le\theta\bm x\right\},
\end{equation}
since $\bm x=\bm v$, $\theta=\rho$ is admissible, while $M\bm x\le\theta\bm x$
paired with $\bm w$ gives
$\rho\langle\bm w,\bm x\rangle\le\theta\langle\bm w,\bm x\rangle$ with
$\langle\bm w,\bm x\rangle>0$. Substituting $M=A_1\operatorname{diag}(\bm g)$
through~\eqref{M-split} turns~\eqref{M-t-program} into
\begin{equation}\label{s-gp-eq}
\begin{aligned}
&\min_{\theta,\bm x,\bm t}\ \theta\\
&\mbox{subject to}\quad
\sum_{l=1}^{N}\left(A_1\right)_{pl}
\left\lbrack(1-\varepsilon_1)+\frac{\varepsilon_1}{t_l}\right\rbrack
\frac{x_l}{\theta\,x_p}\le 1,\quad p=1:N,\\
&\phantom{\mbox{subject to}}\quad
\sum_{j=1}^{n}\frac{\sigma_{1,p}}{c_i}\,t_p\le 1,\quad i=1:m,\qquad
\frac{1}{t_p}\le 1,\qquad \bm x>0,\ \theta>0,
\end{aligned}
\end{equation}
a geometric program~\cite{2007Boyd_gp}: every left-hand side is a posynomial
and the objective function $\theta$ a monomial. In the variables $\log\theta$,
$\log\bm x$ and $\log\bm t$ the constraints become sums of exponentials of
affine functions, hence convex. Only inequalities occur, and this matters: a
geometric program admits posynomial inequality constraints but only monomial
equalities, whereas the equalities of~\eqref{M-constraints} are posynomial.

For three doses this construction is not available: $g_p$ in~\eqref{M-g-def} is
then a quotient of posynomials whose denominator has more than one term, so the
constraint $A_1\operatorname{diag}(\bm g)\bm x\le\theta\bm x$ is not a
posynomial inequality.

\bibliographystyle{plain}
\bibliography{bib_mb}

\end{document}
